\scanpage{29}
\deleted{OPS de plus que $\pi_1(X)$ \illegible{1} un centre trivial. \unsure{En} car. $p>0$, notons qu'il \unsure{suffit} de \unsure{travailler} \unsure{avec} les groupes \unsure{fondamentaux} \unsure{premiers} à $p$ -- et $\unsure{\overline m|=}$ sur les \unsure{localisations|localisés} en $\ell\ne p$ des \unsure{groupes} \unsure{fondamentaux}.}
\mspage{14}

La situation $X,Y,f,g$ provient, par extension du corps de base, d'une situation analogue sur un corps $K$ \inserted{extension} de type fini \deleted{de} $\mathbb Q$. \deleted{\unsure{Choisissons} une} \inserted{Soit $\overline K$ la} clôture algébrique \deleted{$\unsure{\overline K|\overline Q}$} de $K$ dans $k$, \deleted{\unsure{des $k$ l'|des $f_k\,\ell$}} \inserted{et $\overline X,\overline Y,\overline f,\overline g$ \deleted{\unsure{sur|se}}} déduite de \inserted{$X,Y$,} $f,g$ par extension du corps de base de $K$ à $\overline K$. On a donc $\pi_1(\overline f)\ne\pi_1(\overline g)$ satisfaisant la condition d) avec de plus $\pi_1(\overline X)$ a centre trivial.

Mais ces \unsure{hypothèses|hyp. mises} impliquent que les extensions $\mathcal E_{\inserted{X/K}}\deleted{(X)}\inserted{=\pi_1(X)}$, $\mathcal E_{\inserted{Y/K}}\deleted{(Y)}\inserted{=\pi_1(Y)}$ de $\mathcal E_{\overline K/K}$ par $\pi_1(\overline X)$, resp. $\pi_1(\overline Y)$, ainsi que les hom $\pi_1(f)$, $\pi_1(g)$ induits, sont réconstruits à partir de \inserted{$\pi_1(\overline X)$, $\pi_1(\overline Y)$, de} $\pi_1(\overline f)$, $\pi_1(\overline g)$ et de l'action extérieure de $\mathcal E_{\overline K/K}$ sur ces groupes. On va montrer maintenant le
\margin{28}
